\documentclass[11pt]{article}
\usepackage[a4paper,margin=1in]{geometry}
\usepackage{amsmath,amssymb,amsthm}
\usepackage{booktabs}
\usepackage{xcolor}
\usepackage{hyperref}
\hypersetup{colorlinks=true,linkcolor=blue,urlcolor=blue,citecolor=blue}
\newcommand{\fn}[1]{\texttt{#1}}
\newcommand{\A}{\mathcal{A}}
\newcommand{\C}{\mathcal{C}}
\newcommand{\HH}{\mathcal{H}}
\theoremstyle{plain}
\newtheorem{proposition}{Proposition}
\newtheorem{theorem}{Theorem}
\theoremstyle{remark}
\newtheorem{remark}{Remark}

\title{\textbf{Report R26 --- Temperley--Lieb and the four unitary rules}\\
\large What maps, what does not, and why the two families scale differently}
\author{Tier 1 analysis}
\date{2026-08-22}

\begin{document}
\maketitle
\sloppy
\emergencystretch=3em

\begin{abstract}
The four unitary single-$V$ rules $W_{108}, W_{201}, W_{156}, W_{198}$ have been
described in this project as mapping to a Temperley--Lieb (TL) model. This
report establishes exactly how much of that is true. \textbf{One:} the gate is
$g_i = \mathbb{1} - 2e_i$ with $e_i$ an exact orthogonal projector, and on the
full $2^N$ space the TL relation fails for all four --- what holds instead,
exactly, is $e_i e_{i\pm1} e_i = \tau\, e_i P^{\rm spec}$ with a spectator
projector two sites away. \textbf{Two:} on the \emph{kink module} of a wall-free
gap the spectator is pinned and the obstruction disappears: for $W_{156}/W_{198}$
the generators become rank one, the relations hold exactly with arbitrarily many
generators, and the algebra is genuinely $\mathrm{TL}$ at
$\delta = 2\sqrt2$, with gap algebra $\mathbb{C}\oplus M_g(\mathbb{C})$ of
dimension $g^2+1$. \textbf{Three:} $W_{201}/W_{108}$ are \emph{not} TL on any
space we can find; their gap module is the hard-core chain and their generators
have rank growing with the gap. \textbf{Four:} that single structural difference
--- gap module linear in the gap length versus exponential --- is the entire
origin of the different $n_{\rm wcc}$ and $d_{\max}$ scalings, with no
singularity at any TL parameter required. \textbf{Five:} the fragmentation is
quantum, and we exhibit the state that makes it so: a one-dimensional dark line
on each gap, the Rokhsar--Kivelson superposition over domain-wall positions with
amplitude ratio exactly $\sqrt2-1 = \tan(\pi/8)$, which is not a product state
in the site basis. \textbf{Six:} the TL--fragmentation literature works at
$\delta = m$, an \emph{integer} set by a colour count; our $\delta = 2\sqrt2$ is
irrational and arises on a \emph{qubit} chain from a Hadamard, so the model is
outside the family that has been studied. Finally we record that two premises of
the question as posed do not survive: no linear MPO growth has ever been
measured for any of these rules, and the bounded-versus-exponential comparison
that motivated the question was between dissipative children, not between these
unitary parents.
\end{abstract}

\tableofcontents

\section{Temperley--Lieb in one page, for a physicist who has not met it}
\label{sec:tl}

The Temperley--Lieb algebra $\mathrm{TL}_n(\delta)$ is generated by
$e_1,\dots,e_{n}$ subject to three relations,
\begin{equation}
  e_i^2 = \delta\, e_i, \qquad
  e_i e_{i\pm1} e_i = e_i, \qquad
  [e_i, e_j] = 0 \;\; (|i-j| \ge 2).
  \label{eq:tl}
\end{equation}
The single parameter $\delta$ is the \emph{loop parameter}: in the diagrammatic
picture the $e_i$ are cup--cap diagrams, multiplication is stacking, and every
closed loop produced may be erased at the cost of a factor $\delta$. That is the
whole content of the first relation.

Two normalisations circulate and it is worth fixing ours. If one prefers the
generators to be \emph{projectors}, put $e_i = E_i/\delta$; then
\begin{equation}
  e_i^2 = e_i, \qquad e_i e_{i\pm1} e_i = \tau\, e_i, \qquad \tau = \delta^{-2}.
  \label{eq:tlproj}
\end{equation}
We use \eqref{eq:tlproj} throughout, because the objects that arise from our
gates are projectors on the nose. So $\delta = 1/\sqrt\tau$.

Three facts a newcomer needs.

\paragraph{$\delta$ decides everything.} $\mathrm{TL}_n(\delta)$ is semisimple
for generic $\delta$, and its representation theory degenerates at
$\delta = 2\cos(\pi/(m+1))$. The value $\delta = 2$ is the boundary.

\paragraph{Not every $\delta$ is reachable by projectors.} If the $e_i$ are
\emph{orthogonal} projectors on a Hilbert space and \eqref{eq:tlproj} holds for
arbitrarily many generators, then Jones' index theorem forces
$\delta^2 \in \{4\cos^2(\pi/m)\}_{m\ge3} \cup [4,\infty)$. A $\delta$ strictly
between the discrete values and below $2$ is impossible. We use this below to
rule out one rule outright.

\paragraph{Why it appears in fragmentation.} If the local terms of a Hamiltonian
or circuit generate $\mathrm{TL}$, then the conserved quantities are the
commutant of $\mathrm{TL}$, which is known in closed form. When that commutant
grows exponentially in system size one has Hilbert-space fragmentation, and when
its sectors are \emph{not} spanned by product states the fragmentation is called
quantum~\cite{MM2022, Li2023}.

\section{The generators, and the obstruction on the full chain}
\label{sec:obstruction}

Each of the four rules applies a Hadamard at site $i$ conditioned on the two
neighbours being in one fixed pattern $(m^*, n^*)$, the \emph{$V$-context}:
\[
  W_{201}: (0,0), \quad W_{108}: (1,1), \quad
  W_{156}: (1,0), \quad W_{198}: (0,1).
\]
Since $H$ is a reflection, $H = \mathbb{1} - 2|h_-\rangle\langle h_-|$ with
$|h_-\rangle$ its $-1$ eigenvector, the controlled gate is
\begin{equation}
  g_i = \mathbb{1} - 2 e_i, \qquad
  e_i = P^{m^*}_{i-1} \otimes |h_-\rangle\langle h_-|_i \otimes P^{n^*}_{i+1},
  \label{eq:gen}
\end{equation}
and $e_i$ is a product of commuting projectors on disjoint sites, hence an exact
orthogonal projector. Measured: $\lVert e_i^2 - e_i\rVert = 0$ exactly, and
$[e_i,e_j] = 0$ exactly for $|i-j| \ge 2$, all four rules, $N = 5..10$. So two of
the three TL relations \eqref{eq:tlproj} hold identically.

The third does not.

\begin{proposition}[The spectator identity]
\label{prop:spec}
For all four rules, exactly,
\[
  e_i\, e_{i+1}\, e_i = \tau\; e_i\, P^{n^*}_{i+2},
  \qquad
  e_i\, e_{i-1}\, e_i = \tau\; e_i\, P^{m^*}_{i-2},
  \qquad
  \tau = |\langle m^*|h_-\rangle|^2\,|\langle n^*|h_-\rangle|^2 .
\]
\end{proposition}

\noindent Verified to $\le 1.4\times10^{-17}$ for $N = 5..10$. The middle
generator drags in a projector \emph{two sites beyond} the window of $e_i$,
which $e_i$ itself does not carry. On the full $2^N$ space that spectator is a
genuine obstruction and the TL relation fails --- residuals $1.07\times10^{-1}$
($W_{156}/W_{198}$), $1.83\times10^{-2}$ ($W_{201}$), $6.22\times10^{-1}$
($W_{108}$). It fails on the hard-core constrained sector too.

With $|\langle 0|h_-\rangle|^2 = 1/(4+2\sqrt2)$ and
$|\langle 1|h_-\rangle|^2 = (3+2\sqrt2)/(4+2\sqrt2)$ the coefficient is
\begin{center}
\begin{tabular}{llll}
\toprule
rule & $V$-context & $\tau$ & $1/\sqrt\tau$ \\
\midrule
$W_{108}$ & $(1,1)$ symmetric & $0.728553$ & $4-2\sqrt2 = 1.171573$ \\
$W_{156}, W_{198}$ & $(1,0), (0,1)$ domain wall & $1/8$ exactly & $2\sqrt2 = 2.828427$ \\
$W_{201}$ & $(0,0)$ symmetric & $0.021447$ & $4+2\sqrt2 = 6.828427$ \\
\bottomrule
\end{tabular}
\end{center}
The asymmetric case picks up one large and one small overlap, and their product
is $1/8$ exactly; the symmetric cases square a single overlap. Note
$(4-2\sqrt2)(4+2\sqrt2) = 8 = (2\sqrt2)^2$.

\begin{remark}
Only the middle row is a TL parameter. For $W_{108}/W_{201}$ the number
$4\mp2\sqrt2$ is the coefficient in Proposition~\ref{prop:spec} and nothing more,
because \S\ref{sec:notTL} shows those rules never become TL. It should not be
printed in a ``$\delta$'' column, which would imply a structure that is absent.
\end{remark}

\section{Where the obstruction disappears: the kink module}
\label{sec:kink}

The spectator identity says TL holds wherever the far site is \emph{already}
pinned to the $V$-context value. That is exactly what happens inside a wall-free
gap, and it is why the correct variables are not the site-local ones.

For $W_{156}$ the $V$-context is $(1,0)$ and a gap admits no ascent, so its basis
states are the descending configurations $1^k 0^{\,l-k}$. Write $|k\rangle$ for
the position of the domain wall. Then $g_i$ is active exactly on
$\{|i\rangle, |i+1\rangle\}$ and $e_i$ becomes rank one.

\begin{proposition}[TL on the kink chain]
\label{prop:kink}
On the kink module of a wall-free gap of $W_{156}$ or $W_{198}$, each live
generator is $e_i = |v_i\rangle\langle v_i|$ with
$v_i \propto |i\rangle - (1+\sqrt2)\,|i+1\rangle$, adjacent overlaps
$|\langle v_i|v_{i+1}\rangle|^2 = 1/8$ exactly, and the relations
\eqref{eq:tlproj} hold with $\tau = 1/8$, i.e.\ $\delta = 2\sqrt2$, for
\emph{every} adjacent pair, with the number of generators growing as the gap.
\end{proposition}

\noindent Verified independently here, by building the module from the site
operators \eqref{eq:gen} and projecting --- not by re-implementing the algebra:

\begin{center}
\begin{tabular}{rrrrrrrr}
\toprule
$N$ & module dim & live gens & ranks & closure & TL error & $|\langle v_i|v_{i+1}\rangle|^2$ & ratio \\
\midrule
5  & 6  & 3 & $\{1\}$ & $0$ & $1.4\times10^{-17}$ & $0.125000$ & $-2.414214$ \\
8  & 9  & 6 & $\{1\}$ & $0$ & $1.4\times10^{-17}$ & $0.125000$ & $-2.414214$ \\
10 & 11 & 8 & $\{1\}$ & $0$ & $1.4\times10^{-17}$ & $0.125000$ & $-2.414214$ \\
\bottomrule
\end{tabular}
\end{center}

\noindent with $-2.414214 = -(1+\sqrt2)$ and $0.125 = 1/8$ exactly, and distant
generators commuting to $0$. The gap algebra is
\begin{equation}
  \A_{\rm gap} \;\cong\; \mathbb{C} \oplus M_g(\mathbb{C}),
  \qquad \dim \A_{\rm gap} = g^2 + 1,
\end{equation}
measured as $10, 17, 26, 37, 50, 65$ for $g = 3,\dots,8$. Since
$\delta = 2\sqrt2 > 2$ this is the standard semisimple-TL statement, so TL is not
decoration here --- it is what proves the gap decomposition.

\begin{remark}[A trap worth recording]
Testing \eqref{eq:tlproj} in the site-local variables on $2^N$ gives ``longest
TL run $= 2$'' for all four rules, which reads as a decisive negative. It is an
artefact: for $W_{156}/W_{198}$ three adjacent generators are never
simultaneously live in any sector, so $\mathrm{TL}_3$ is never even testable
there. The whole chain is a \emph{product over gaps} of TL algebras with the gap
boundaries frozen, not one TL algebra on $2^N$. This report's first draft drew
the wrong conclusion from exactly this.
\end{remark}

\section{\texorpdfstring{$W_{201}$ and $W_{108}$ are not Temperley--Lieb}%
{W201 and W108 are not Temperley-Lieb}}
\label{sec:notTL}

For the symmetric rules the wall word is $11$ or $00$, so a wall-free gap is not
a kink chain but the \emph{hard-core} chain --- configurations with no two
adjacent $1$s (resp.\ $0$s) --- whose dimension is exponential in the gap length.
On their own largest block the generators are nowhere near rank one: at $N=10$,
$W_{201}$ at obc0 has module dimension $144$ with generator ranks
$\{34,39,40,42,55\}$, and $W_{108}$ module dimension $55$ with ranks
$\{13,15,16,21\}$. The TL residuals stay at $1.8\times10^{-2}$ and
$6.2\times10^{-1}$.

There are therefore two independent obstructions, and they are different for the
two rules:
\begin{itemize}
\item $W_{108}$: Jones' index. $\delta^2 = (4-2\sqrt2)^2 = 24-16\sqrt2 = 1.3726$
lies in $(1,2)$, and the nearest admissible discrete value $4\cos^2(\pi/3) = 1$
is far away. No representation by orthogonal projectors exists for arbitrarily
many generators, whatever variables are chosen.
\item $W_{201}$: \emph{not} Jones --- $\delta = 6.83 > 2$ is perfectly
admissible. The obstruction is the direct one: the generators have rank growing
with the gap and simply do not satisfy \eqref{eq:tlproj}.
\end{itemize}

\section{Why the two families scale differently}
\label{sec:scaling}

This answers the question that motivated the report, and the answer needs no
singularity.

The minimal wall word is the bit complement of the $V$-context. That one fact
splits the family:
\begin{center}
\begin{tabular}{lllll}
\toprule
rules & $V$-context & wall word & gap module & dimension in gap length $l$ \\
\midrule
$W_{156}, W_{198}$ & $(1,0),(0,1)$ & $01, 10$ & kink chain $1^a0^b$ & $l+1$, \textbf{linear} \\
$W_{201}, W_{108}$ & $(0,0),(1,1)$ & $11, 00$ & hard-core chain & \textbf{exponential} \\
\bottomrule
\end{tabular}
\end{center}

\noindent A gap that admits no ascent has exactly one degree of freedom, the
position of its domain wall. A gap with a hard-core constraint has a Fibonacci
number of states. Largest block dimension, $N = 4..12$: $W_{201}$ at obc0 gives
$8, 13, 21, 34, 55, 89, 144, 233, 377$ (Fibonacci), $W_{108}$ gives
$3,5,8,13,21,34,55,89,144$, both families at pbc give Lucas numbers (hard core on
a ring), while $W_{156}/W_{198}$ give $4,5,6,9,12,16,20,27,36$, products of kink
chains. This is consistent with the $F(N+2)$ found independently in R24 for the
constrained space of $W_{201}$.

Linear versus exponential gap-module dimension is the whole of it: it fixes
$n_{\rm wcc}$, it fixes $d_{\max}$, and it fixes the sector-size statistics.
There is nothing to reconcile \emph{inside} TL because only one family is in TL.
And it answers the entanglement question in the same breath --- a module of
dimension $l+1$ cannot carry what a Fibonacci module can.

\section{The fragmentation is quantum, and here is the state}
\label{sec:quantum}

Fragmentation is called \emph{quantum} when the sectors are not spanned by an
orthonormal product basis. For $W_{156}/W_{198}$ the witness is explicit and we
exhibit it.

The summand $\mathbb{C}$ in $\A_{\rm gap} \cong \mathbb{C}\oplus M_g(\mathbb{C})$
is a \emph{dark line}: a state annihilated by every generator of the gap. Solving
for the common kernel:

\begin{center}
\begin{tabular}{rrrrr}
\toprule
rule & $N$ & kernel dim & amplitude ratio & support (kink states) \\
\midrule
$W_{156}$ & $5,\dots,10$ & $1$ & $0.414214$ & $N-1$ \\
$W_{198}$ & $5,\dots,10$ & $1$ & $2.414214$ & $N-1$ \\
\bottomrule
\end{tabular}
\end{center}

\noindent with $0.414214 = \sqrt2 - 1 = \tan(\pi/8)$ exactly and $2.414214$ its
reciprocal, the mirror rule giving the mirror ratio. So the dark state is the
Rokhsar--Kivelson superposition
\begin{equation}
  |{\rm dark}\rangle \;\propto\; \sum_k t^k |k\rangle,
  \qquad t = \sqrt2 - 1 = \tan(\pi/8),
\end{equation}
over \emph{all} domain-wall positions in the gap. Its support is every kink state
of the gap, so in the computational (site) basis it is manifestly not a product
state, and it cannot be rotated to one without leaving the sector.

\textbf{In which basis, explicitly.} The fragmentation is classical in the kink
coordinate $|k\rangle$ --- the sectors are gaps, labelled by wall positions ---
and quantum in the site basis, because one conserved direction per gap is the
entangled superposition above. Equivalently: the sector decomposition is spanned
by product states, but admits an \emph{orthonormal} product basis only when every
block is a subcube, which a census over the four rules, two boundary conditions
and $N = 3..16$ found in exactly three cases, all at $N = 3$.

\section{The literature: this model is not the one that has been studied}
\label{sec:lit}

The TL--fragmentation literature works at integer $\delta$ on multi-level chains.
Moudgalya and Motrunich~\cite{MM2022} define the TL models on an $m$-level chain
by $\hat e_{j,j+1} = \sum_{\alpha,\beta=1}^{m} (|\alpha\alpha\rangle\langle\beta\beta|)_{j,j+1}$
and state the relations with
\[
  q + q^{-1} = m, \qquad\text{i.e.}\qquad \delta = m \in \{2,3,4,\dots\}.
\]
Li, Sala and Pollmann~\cite{Li2023}, whose open-system TL model is the closest
published neighbour to this project, use the spin-$1$ case $m = 3$: their
$e_{j,j+1}$ projects onto the dimer
$\tfrac{1}{\sqrt3}(|{+}{+}\rangle + |00\rangle + |{-}{-}\rangle)$, so
$\delta = 3$. Klocke, Moore and Buchhold~\cite{Klocke2024} use an
$SU(N)$-symmetric TL Hamiltonian, again with $\delta$ set by a colour count.

Ours is a \textbf{qubit} chain --- two levels --- at
$\delta = 2\sqrt2 = 2.828\ldots$, irrational, arising from the rank-one spectral
projector of a Hadamard rather than from any colour multiplicity. On a two-level
chain the standard construction gives $\delta = m = 2$, exactly the critical
Jones value; we sit strictly above it. The model is therefore \emph{not} a member
of the pair-flip family that the fragmentation literature treats, and we have
found no paper reporting a TL fragmentation model at irrational $\delta$ on
qubits, nor one arising as the bond algebra of a Floquet cellular automaton.

Two further recent papers were checked and do not overlap:
Zhou, Yang and Chen~\cite{Zhou2026} (quantum HSF and entangled frozen states,
TL among four models, entanglement $S\sim\sqrt L$) and Han, Hart, Khudorozhkov
and Nandkishore~\cite{Han2026} (a construction protocol for quantum-HSF
Hamiltonians; TL not used).

\begin{remark}[The question as posed cannot be asked]
Part of the motivating question was whether anyone had previously reported the
``linear evolution-operator MPO growth'' of the $W_{156}/W_{198}$ model. That
measurement does not exist. What has been measured, for rule $156$ alone, is an
$N$-independent master curve with increments of exactly $6$ per period in an
early window, a concave approach thereafter, a departure law
$t_{\rm dep} \approx 0.91N - 8.4$, and an exact plateau
$\chi_{\rm sat} = (N^2+15)/4$ --- \emph{quadratic} in $N$, not linear in $t$.
Those are the results worth asking the literature about, and we have found no
report of them. There is no MPO data at all for $W_{198}, W_{201}, W_{108}$, so
the linear-versus-exponential comparison between the two unitary families has
never been run; the bounded-versus-exponential contrast that motivated the
question was between the \emph{dissipative children} $W_{28}/W_{70}$,
$W_{29}/W_{71}$, $W_{157}/W_{199}$ and $W_{73}/W_{109}$, not between these
parents.
\end{remark}

\section{Summary}

\begin{enumerate}
\item $g_i = \mathbb{1} - 2e_i$ with $e_i$ an exact orthogonal projector; two of
the three TL relations hold identically for all four rules.
\item The third fails on $2^N$, and exactly: $e_ie_{i\pm1}e_i = \tau e_i P^{\rm
spec}$. The spectator is the whole obstruction.
\item On a wall-free gap the spectator is pinned. For $W_{156}/W_{198}$ the
generators become rank one and the algebra is genuinely TL at $\delta = 2\sqrt2$,
with $\A_{\rm gap} = \mathbb{C}\oplus M_g(\mathbb{C})$, $\dim = g^2+1$. The chain
object is a product over gaps, not one TL algebra.
\item $W_{201}/W_{108}$ are not TL anywhere: hard-core gap module, high-rank
generators; and $W_{108}$ is excluded outright by Jones' index, $W_{201}$ is not.
\item Linear versus exponential gap-module dimension explains the different
$n_{\rm wcc}$, $d_{\max}$ and sector statistics with no singularity.
\item The fragmentation is quantum; the witness is a one-dimensional dark RK line
per gap with amplitude ratio $\sqrt2-1 = \tan(\pi/8)$.
\item The published TL fragmentation models sit at integer $\delta = m$ on
$m$-level chains; ours is irrational on qubits and is not in that family.
\end{enumerate}

\section*{Provenance}
\addcontentsline{toc}{section}{Provenance}

The kink-variable construction, the gap-module classification and the
$\mathbb{C}\oplus M_g$ decomposition are the Tier-2 session's, from its reports
R-T15 and R-T17. Every numerical claim attributed to them above was reproduced
here independently, from the site operators \eqref{eq:gen} and a projection onto
the module, without using their code: the rank-one property, the $1/8$ overlaps,
the $-(1+\sqrt2)$ ratio, the TL residuals, $\dim\A = g^2+1$, and the dark-state
ratio $\sqrt2-1$. The spectator identity of Proposition~\ref{prop:spec} and the
three $\tau$ values are this report's. The MPO facts in \S\ref{sec:lit} are the
Schmidt-rank session's and are quoted, not reproduced.

Two errors made in the course of this work are recorded because they are
instructive. A site-ordering inconsistency (operators indexed from the most
significant bit, sector search from the least) was invisible because reflection
maps $W_{156}\leftrightarrow W_{198}$ and fixes $W_{201}, W_{108}$, so each rule
was silently tested against its mirror's sectors. And the Jones argument was
first stated as excluding both symmetric rules; it excludes only $W_{108}$.

\begin{thebibliography}{9}

\bibitem{MM2022} S.~Moudgalya and O.~I.~Motrunich,
  \emph{Hilbert space fragmentation and commutant algebras},
  Phys.\ Rev.\ X \textbf{12}, 011050 (2022); arXiv:2108.10324.

\bibitem{Li2023} Y.~Li, P.~Sala and F.~Pollmann,
  \emph{Hilbert space fragmentation in open quantum systems},
  Phys.\ Rev.\ Research \textbf{5}, 043239 (2023); arXiv:2305.06918.

\bibitem{Klocke2024} K.~Klocke, J.~E.~Moore and M.~Buchhold,
  \emph{Power-law entanglement and Hilbert space fragmentation in non-reciprocal
  quantum circuits},
  Phys.\ Rev.\ Lett.\ \textbf{133}, 070401 (2024); arXiv:2405.06021.

\bibitem{Zhou2026} Z.~Zhou, T.-H.~Yang and B.-T.~Chen,
  \emph{Quantum Hilbert space fragmentation and entangled frozen states},
  arXiv:2604.05218 (2026).

\bibitem{Han2026} Y.~Han, O.~Hart, A.~Khudorozhkov and R.~Nandkishore,
  \emph{Quantum fragmentation},
  arXiv:2604.06461 (2026).

\bibitem{Jones1983} V.~F.~R.~Jones,
  \emph{Index for subfactors},
  Invent.\ Math.\ \textbf{72}, 1 (1983).

\end{thebibliography}

\end{document}
